% =============================================================================
% Chapter 5: Distributed Training Systems — ML Systems Lecture Slides (Vol II)
% =============================================================================
% IMAGES NEEDED (SVGs converted to PDF):
%   - 3d-parallelism-cube.pdf   - data-parallel-flow.pdf
%   - scaling-tax.pdf           - zero-memory.pdf
%   - pipeline-schedule.pdf     - tensor-parallelism.pdf
%   - hybrid-3d.pdf
% =============================================================================
\documentclass[aspectratio=169, 12pt]{beamer}
\usepackage{../../assets/beamerthememlsys}

\mlsyssetup{
  volume       = {Volume II},
  chapter      = {Chapter 5},
  logo         = {../../assets/img/logo-mlsysbook.png},
  instlogo     = {../../assets/img/logo-harvard.png},
  chaptertitle = {Distributed Training Systems},
}

% --- Fonts ---
\usepackage{fontspec}
\setsansfont{Helvetica Neue}[
  BoldFont={Helvetica Neue Bold},
  ItalicFont={Helvetica Neue Italic},
  BoldItalicFont={Helvetica Neue Bold Italic},
]
\setmonofont{JetBrains Mono}[Scale=0.85]

% --- Packages ---
\usepackage{booktabs}
\usepackage{amsmath}

% --- Image paths ---
\graphicspath{
  {images/}
}

% --- Chapter-specific macros ---
\newcommand{\CCR}{Communication-Computation Ratio}

% --- Helper: safe image include ---
\newcommand{\safeimg}[2][width=\textwidth,keepaspectratio]{%
  \IfFileExists{images/#2}{\includegraphics[#1]{#2}}{%
    \IfFileExists{#2}{\includegraphics[#1]{#2}}{%
      \fbox{\parbox[c][2.5cm][c]{0.85\linewidth}{%
        \centering\footnotesize\textcolor{midgray}{[Missing image]}}}%
    }%
  }%
}

% --- Section count for navigation ---
\setcounter{mlsystotalsections}{8}

\title{Distributed Training Systems}
\author{Vijay Janapa Reddi}
\institute{Harvard University}
\date{}

\begin{document}

% =============================================================================
% TITLE SLIDE
% =============================================================================
\mlsystitle{Distributed Training}{The Physics of Scaling}{cover_distributed.png}

% =============================================================================
% LEARNING OBJECTIVES
% =============================================================================
\begin{frame}{Learning Objectives}
\note{[2 min] Read through objectives. Emphasize that this chapter bridges the
physical fleet (Part I) with algorithmic scaling. Ask: ``How many GPUs did it take
to train GPT-3?'' (answer: $\sim$10,000 V100s for weeks).}

\small
\begin{enumerate}
  \item Apply the \textbf{Law of Distributed Efficiency} to diagnose compute-bound vs.\ communication-bound workloads
  \item Implement \textbf{Data Parallelism} with gradient synchronization and AllReduce
  \item Apply \textbf{ZeRO/FSDP} memory sharding and quantify the communication trade-off
  \item Design \textbf{Tensor Parallelism} strategies that exploit intra-node NVLink
  \item Construct \textbf{Pipeline Parallelism} schedules to minimize bubble overhead
  \item Select \textbf{Synchronization Models} (BSP, SSP, Async) based on cluster constraints
  \item Architect \textbf{Hybrid 3D Parallelism} configurations for frontier models
\end{enumerate}

\end{frame}

\begin{frame}{Visual Language}
\note{[1 min] Explain the semantic color system used throughout the course.
These colors are consistent across all diagrams and slides.}

\small
Throughout this course, colors carry meaning:

\vspace{0.3cm}
\begin{columns}[T]
  \begin{column}{0.45\textwidth}
    \begin{mlsyscard}{computestroke}
    \textbf{Blue} --- Compute / Processing\\
    {\footnotesize GPU ops, forward/backward pass, inference}
    \end{mlsyscard}
    \vspace{0.15cm}
    \begin{mlsyscard}{datastroke}
    \textbf{Green} --- Data / Memory\\
    {\footnotesize Data flow, caches, healthy paths}
    \end{mlsyscard}
  \end{column}
  \begin{column}{0.45\textwidth}
    \begin{mlsyscard}{routingstroke}
    \textbf{Orange} --- Routing / Scheduling\\
    {\footnotesize Load balancers, batch windows}
    \end{mlsyscard}
    \vspace{0.15cm}
    \begin{mlsyscard}{errorstroke}
    \textbf{Red} --- Error / Cost / Bottleneck\\
    {\footnotesize Loss, decode phase, waste}
    \end{mlsyscard}
  \end{column}
\end{columns}
\end{frame}



% =============================================================================
\section{Why Distribute?}
% =============================================================================

\begin{frame}{The Single-GPU Wall}
\note{[3 min] Open with the impossibility argument. GPT-3 on one V100 = 355 years.
The physics forces distribution. Ask: ``What are the three signals that force you
to go distributed?'' Memory, time, and data scale.}

\small
\begin{columns}[T]
  \begin{column}{0.55\textwidth}
    A single GPU cannot train frontier models:

    \vspace{0.15cm}
    \begin{mlsyscard}{errorstroke}
    \textbf{Memory}: GPT-3 (175B) needs 350 GB in BF16\\
    {\footnotesize A single H100 has only 80 GB HBM3}
    \end{mlsyscard}

    \vspace{0.1cm}
    \begin{mlsyscard}{errorstroke}
    \textbf{Time}: GPT-3 on a single V100 $\approx$ \textbf{355 years}
    \end{mlsyscard}

    \vspace{0.1cm}
    \begin{mlsyscard}{errorstroke}
    \textbf{Data}: Training data reaches multiple TB\\
    {\footnotesize Exceeds single-machine storage}
    \end{mlsyscard}
  \end{column}
  \begin{column}{0.42\textwidth}
    \vspace{0.2cm}
    \mlsysconcept{Key insight:}{Distribution is not optional --- it is a \emph{physical necessity}.}

    \vspace{0.15cm}
    \footnotesize
    The art is managing the \textbf{Coordination Tax}: synchronization overhead that grows with scale.
  \end{column}
\end{columns}

\end{frame}

\begin{frame}{The Iron Law of Scale}
\note{[3 min] Connect to the Fleet Stack from Ch1. The step time equation governs
everything. Walk through each term. The CCR determines whether you have a
supercomputer or a collection of idling heaters.}

\small
$$T_{\text{step}}(N) = \frac{T_{\text{compute}}}{N} + T_{\text{comm}}(N) - T_{\text{overlap}}$$

\vspace{0.1cm}
\begin{columns}[T]
  \begin{column}{0.30\textwidth}
    \centering
    \textcolor{computestroke}{\textbf{Compute}}\\[0.05cm]
    {\scriptsize $T_{\text{compute}}/N$\\Scales perfectly\\with GPUs}
  \end{column}
  \begin{column}{0.30\textwidth}
    \centering
    \textcolor{errorstroke}{\textbf{Communication}}\\[0.05cm]
    {\scriptsize $T_{\text{comm}}(N)$\\Grows with scale\\(gradient sync)}
  \end{column}
  \begin{column}{0.30\textwidth}
    \centering
    \textcolor{datastroke}{\textbf{Overlap}}\\[0.05cm]
    {\scriptsize $T_{\text{overlap}}$\\Comm hidden behind\\computation}
  \end{column}
\end{columns}

\vspace{0.2cm}
\begin{mlsyscard}{crimson}
The \textbf{\CCR{}} ($T_{\text{comm}}/T_{\text{compute}}$) determines whether your cluster is a supercomputer or a collection of idling heaters.
\end{mlsyscard}

\end{frame}

\begin{frame}{The 3D Parallelism Cube}
\note{[3 min] This is the conceptual framework for the entire chapter. Three
orthogonal axes, each independent. Total GPUs = d $\times$ t $\times$ p.
Ask: ``Which axis would you use first and why?''}

% --- Full-width image ---
\centering
\safeimg[width=0.88\textwidth,height=4.5cm,keepaspectratio]{3d-parallelism-cube.pdf}

\vspace{0.1cm}
\footnotesize
\begin{columns}[T]
  \begin{column}{0.30\textwidth}
    \centering
    \textcolor{computestroke}{\textbf{Data (d)}}\\
    Replicate model, split data
  \end{column}
  \begin{column}{0.30\textwidth}
    \centering
    \textcolor{routingstroke}{\textbf{Tensor (t)}}\\
    Split ops within layers
  \end{column}
  \begin{column}{0.30\textwidth}
    \centering
    \textcolor{datastroke}{\textbf{Pipeline (p)}}\\
    Split layers across stages
  \end{column}
\end{columns}

\end{frame}

% =============================================================================
\section{Data Parallelism}
% =============================================================================

\begin{frame}{Data Parallelism: The Default Strategy}
\note{[3 min] Simplest approach: replicate the model, split the data. Walk through
the five stages. Key insight: gradient averaging is mathematically equivalent to
single-GPU training with a larger batch. The bottleneck is AllReduce.}

% --- Full-width image ---
\centering
\safeimg[width=0.92\textwidth,height=4.8cm,keepaspectratio]{data-parallel-flow.pdf}

\vspace{0.1cm}
{\small The \textbf{critical synchronization point} is Stage 4: AllReduce must complete before any GPU updates parameters.}

\end{frame}

\begin{frame}{Data Parallelism: The Math}
\note{[2 min] Walk through the gradient equivalence proof. This is why DP works:
linearity of expectation. Each GPU computes gradients independently, and averaging
them is the same as computing on the combined batch.}

\footnotesize
\textbf{Gradient averaging preserves SGD equivalence:}

\vspace{0.05cm}
Each GPU $k$ computes local gradients on batch $B_k$:
$$g_k = \frac{1}{|B_k|} \sum_{x_i \in B_k} \nabla_\theta L(\theta, x_i)$$

\vspace{-0.1cm}
Global update averages local gradients:
$$g_{\text{global}} = \frac{1}{N} \sum_{k=1}^N g_k = \frac{1}{|B_{\text{total}}|} \sum_{x_i \in B_{\text{total}}} \nabla_\theta L(\theta, x_i)$$

\vspace{-0.05cm}
\mlsysconcept{Key:}{Mathematically equivalent to single-device training with batch size $N \times B$.}

\end{frame}

% --- ACTIVE LEARNING 1: Predict ---
\begin{frame}{Predict: Where Does DP Break?}
\note{[2 min] Prediction exercise. Give students 60 seconds to identify the three
ceilings. Do NOT reveal yet. This primes the Communication Wall discussion.}

\centering
\vspace{1.0cm}
{\Large\bfseries Think--Write--Share}

\vspace{0.5cm}
{\large Data parallelism replicates the entire model on each GPU.\\
What are the \alert{three hard ceilings} that limit this strategy?}

\vspace{0.5cm}
{\normalsize Write down 3 failure modes. \textcolor{midgray}{(60 seconds)}}

\vspace{0.5cm}
{\small\textcolor{lightgray}{Turn to a neighbor and compare your answers.}}

\end{frame}

\begin{frame}{The Three Ceilings of Data Parallelism}
\note{[3 min] Reveal after prediction. Memory wall: 175B model needs 1+ TB per GPU.
Bandwidth wall: AllReduce dominates at scale. Batch size trap: critical batch size
means diminishing returns. These three ceilings motivate the rest of the chapter.}

\footnotesize
\renewcommand{\arraystretch}{1.15}
\begin{tabular}{@{}lp{3.5cm}p{5.5cm}@{}}
  \toprule
  \textbf{Ceiling} & \textbf{Mechanism} & \textbf{Quantitative Evidence} \\
  \midrule
  \textcolor{errorstroke}{\textbf{Memory Wall}} & Every GPU holds full model + optimizer & 175B model: 1+ TB/GPU (16 bytes/param with Adam) \\
  \textcolor{errorstroke}{\textbf{Bandwidth Wall}} & AllReduce cost $\propto 2(N\!-\!1)/N \times M$ & LLM gradient sync: 50\%+ of step time at 1,024 GPUs \\
  \textcolor{errorstroke}{\textbf{Batch Size Trap}} & $B_{\text{global}} = N \times B_{\text{local}}$ & Critical batch size: 8K (ResNet), 4M tokens (LLMs) \\
  \bottomrule
\end{tabular}

\vspace{0.15cm}
\begin{mlsyscard}{crimson}
{\footnotesize Beyond the \textbf{Critical Batch Size}, doubling GPUs doubles cost but yields minimal convergence speedup. Data parallelism alone cannot train frontier models.}
\end{mlsyscard}

\end{frame}

% =============================================================================
\section{Scaling Efficiency}
% =============================================================================

\begin{frame}{The Scaling Tax}
\note{[3 min] Walk through the scaling tax figure. Compute-bound models (ResNet)
scale well. Bandwidth-bound models (LLMs) hit a communication wall.
The gap between ideal and actual is the ``tax'' you pay for coordination.}

% --- Full-width image ---
\centering
\safeimg[width=0.88\textwidth,height=4.5cm,keepaspectratio]{scaling-tax.pdf}

\vspace{0.1cm}
{\small $\text{Speedup}(N) = \frac{N}{1 + (N-1) \cdot r}$ where $r$ = communication fraction of step time.}

\end{frame}

\mlsysfocus{Scaling Efficiency Bound}{%
$\text{Efficiency}(N) = \dfrac{1}{1 + \dfrac{N(T_{\text{comm}} - T_{\text{overlap}})}{T_{\text{compute}}}}$%
}

\begin{frame}{Synchronization Models: BSP, SSP, Async}
\note{[3 min] Three choices for when workers synchronize. BSP is the default
(mathematically equivalent to single-GPU). SSP allows bounded staleness.
Async maximizes throughput but degrades convergence. The trade-off is
consistency vs.\ throughput.}

\footnotesize
\renewcommand{\arraystretch}{1.15}
\begin{tabular}{@{}lp{1.8cm}p{2cm}p{2.5cm}p{3cm}@{}}
  \toprule
  \textbf{Model} & \textbf{Staleness} & \textbf{Throughput} & \textbf{Convergence} & \textbf{Use Case} \\
  \midrule
  \textbf{BSP}   & $\tau = 0$       & Slowest worker & Equivalent to 1-GPU & Final training runs \\
  \textbf{SSP}   & $\tau \leq s$    & Higher than BSP & Near-equivalent      & Hyperparameter search \\
  \textbf{Async}  & $\tau = \infty$  & Maximum         & Degraded (15--30\%)  & Large heterogeneous clusters \\
  \bottomrule
\end{tabular}

\vspace{0.15cm}
\begin{columns}[T]
  \begin{column}{0.48\textwidth}
    \begin{mlsyscard}{computestroke}
    \textbf{BSP}: All workers barrier every step.\\
    {\scriptsize Straggler problem: 1 slow GPU delays all 1,024.}
    \end{mlsyscard}
  \end{column}
  \begin{column}{0.48\textwidth}
    \begin{mlsyscard}{routingstroke}
    \textbf{Compensation}: $\eta' = \eta / \sqrt{1 + \tau}$\\
    {\scriptsize Reduce LR to counteract gradient staleness.}
    \end{mlsyscard}
  \end{column}
\end{columns}

\end{frame}

\begin{frame}{Critical Batch Size: When Does Parallelism Hurt?}
\note{[2 min] The critical batch size is where gradient noise equals signal.
Below it: linear scaling works. Above it: diminishing returns. Three orders
of magnitude variation across tasks. This limits how many GPUs you can usefully add.}

\small
\begin{columns}[T]
  \begin{column}{0.55\textwidth}
    $$B^* \approx \frac{\text{tr}(\Sigma)}{\|\nabla L(\theta)\|^2}$$

    \vspace{0.1cm}
    {\footnotesize Batch size where gradient noise = signal strength.}

    \vspace{0.15cm}
    \textbf{Empirical values:}
    \begin{itemize}\setlength\itemsep{0pt}
      \item {\footnotesize ResNet-50: $B^* \approx$ 8K--16K}
      \item {\footnotesize BERT-Large: $B^* \approx$ 32K--65K}
      \item {\footnotesize GPT-3 scale: $B^* \approx$ 1M--4M tokens}
    \end{itemize}
  \end{column}
  \begin{column}{0.42\textwidth}
    \begin{mlsyscard}{errorstroke}
    \textbf{Below $B^*$}: Linear scaling.\\
    Doubling batch halves iterations.\\[0.1cm]
    \textbf{Above $B^*$}: Diminishing returns.\\
    More GPUs $\neq$ faster convergence.
    \end{mlsyscard}

    \vspace{0.1cm}
    {\scriptsize Adding workers beyond $B^*/b$ wastes compute without reducing training time.}
  \end{column}
\end{columns}

\end{frame}

% =============================================================================
\section{ZeRO \& FSDP}
% =============================================================================

\begin{frame}{ZeRO: Progressive Memory Sharding}
\note{[3 min] Walk through the three stages. ZeRO-1 shards optimizer (4x savings),
ZeRO-2 adds gradients (8x), ZeRO-3 adds parameters (Nx linear). The trade-off
is always more communication for less memory. Use the 7B model example:
112 GB baseline drops to 1.75 GB with ZeRO-3 on 64 GPUs.}

% --- Full-width image ---
\centering
\safeimg[width=0.92\textwidth,height=4.8cm,keepaspectratio]{zero-memory.pdf}

\vspace{0.1cm}
{\small \textbf{7B model}: 112 GB (DDP) $\to$ 1.75 GB (ZeRO-3, 64 GPUs). Trade-off: 10--25\% comm overhead.}

\end{frame}

% --- ACTIVE LEARNING: Micro-Retrieval Cue ---
\begin{frame}{Quick Check}
\note{[1 min] Answer: 1120/80 = 14 GPUs minimum. Plus activations, so ~16-20.}

\centering
\vspace{1.0cm}
{\Large\bfseries Quick Check}

\vspace{0.6cm}
{\large A 70B model needs 1,120 GB of training state.\\How many 80 GB GPUs minimum with ZeRO-3?}

\vspace{0.5cm}
{\normalsize\textcolor{midgray}{15 seconds --- then I will cold-call.}}

\end{frame}



% --- ACTIVE LEARNING 2: Exercise ---
\begin{frame}{Your Turn: ZeRO Memory Budget}
\note{[3 min] Give students 90 seconds. Answer: 175B $\times$ 16 bytes = 2,800 GB.
With 64 GPUs: 2,800/64 = 43.75 GB per GPU. Plus 50 GB activations = 93.75 GB.
Exceeds 80 GB. Pure DP fails even with ZeRO-3.}

\footnotesize
\begin{columns}[T]
  \begin{column}{0.55\textwidth}
    {\normalsize\bfseries Memory Wall Calculation}

    \vspace{0.15cm}
    Training a \textbf{175B parameter} model (GPT-3):
    \begin{itemize}\setlength\itemsep{0pt}
      \item Parameters: 2 bytes (FP16) each
      \item Optimizer: 12 bytes (Adam FP32) each
      \item Total: 16 bytes/parameter
      \item Activations: $\sim$50 GB per GPU
      \item Hardware: A100 (80 GB HBM)
    \end{itemize}

    \vspace{0.1cm}
    \textbf{With ZeRO-3 on 64 GPUs, does it fit?}

    {\scriptsize\textcolor{midgray}{(90 seconds --- then compare with a neighbor)}}
  \end{column}
  \begin{column}{0.42\textwidth}
    \pause
    \begin{mlsyscard}{errorstroke}
    \textbf{Solution:}\\[0.05cm]
    {\scriptsize
    Static: $175\text{B} \times 16 = 2{,}800$ GB\\
    Per-GPU: $2{,}800 / 64 \approx 44$ GB\\
    + Activations: $\sim$50 GB\\
    Total: $\sim$94 GB\\[0.05cm]
    \textcolor{errorstroke}{\textbf{Exceeds 80 GB!}}\\
    Pure DP fails. Need \textbf{Tensor Parallelism}.
    }
    \end{mlsyscard}
  \end{column}
\end{columns}

\end{frame}

% =============================================================================
\section{Pipeline Parallelism}
% =============================================================================

\begin{frame}{Pipeline Parallelism: Splitting Depth}
\note{[3 min] Walk through the pipeline schedule. Forward passes cascade diagonally,
backward passes follow in reverse. The gray blocks are ``bubbles'' --- idle GPU time.
The bubble fraction is (P-1)/(P-1+M). With 4 stages, 4 microbatches: 43\% idle.
Ask: ``How do you shrink the bubble?''}

% --- Full-width image ---
\centering
\safeimg[width=0.92\textwidth,height=4.5cm,keepaspectratio]{pipeline-schedule.pdf}

\vspace{0.1cm}
{\small \textbf{Bubble Tax}: $(P-1)/(P-1+M)$. With $P=8, M=32$: only 18\% idle. \alert{Key}: $M \gg P$.}

\end{frame}

\begin{frame}{Pipeline Parallelism: Communication Advantage}
\note{[2 min] The key advantage: only activation tensors cross device boundaries.
For hidden dim 8192 in BF16 with micro-batch 1, that is 16 KB per boundary.
Compare to GB-scale gradient AllReduce in data parallelism. This is why PP
works across nodes on InfiniBand.}

\small
\begin{columns}[T]
  \begin{column}{0.55\textwidth}
    \textbf{Why pipeline parallelism scales across nodes:}

    \vspace{0.15cm}
    {\footnotesize
    \renewcommand{\arraystretch}{1.15}
    \begin{tabular}{@{}lll@{}}
      \toprule
      \textbf{Strategy} & \textbf{Comm.\ Volume} & \textbf{Interconnect} \\
      \midrule
      Data Par. & $\propto M$ (model size) & High BW \\
      Tensor Par. & $\propto B \times L$ & NVLink \\
      Pipeline Par. & $\propto B \times H$ & InfiniBand OK \\
      \bottomrule
    \end{tabular}
    }

    \vspace{0.15cm}
    \mlsysconcept{Advantage:}{Only boundary activations cross devices --- 16 KB vs.\ GB.}
  \end{column}
  \begin{column}{0.42\textwidth}
    \begin{mlsyscard}{errorstroke}
    \textbf{Pipeline Bubble Cost}\\[0.1cm]
    {\footnotesize
    $P=8$ stages, $M=32$ micro-batches:\\
    Bubble = $7/39 \approx$ \textbf{18\%} idle\\[0.1cm]
    $P=16$ stages, $M=16$ micro-batches:\\
    Bubble = $15/31 \approx$ \textbf{48\%} idle!
    }
    \end{mlsyscard}

    \vspace{0.1cm}
    {\scriptsize \textbf{1F1B schedule}: interleave forward/backward to reduce peak memory from $O(M \times P)$ to $O(P)$.}
  \end{column}
\end{columns}

\end{frame}

% =============================================================================
\section{Tensor Parallelism}
% =============================================================================

\begin{frame}{Tensor Parallelism: Splitting Operations}
\note{[3 min] Walk through the Megatron-LM column+row splitting pattern.
Column-parallel splits the first linear, activation is local (no comm needed),
row-parallel computes partial results, then AllReduce combines them. The critical
constraint: AllReduce on EVERY layer requires NVLink bandwidth. Never run TP
across racks.}

% --- Full-width image ---
\centering
\safeimg[width=0.92\textwidth,height=4.8cm,keepaspectratio]{tensor-parallelism.pdf}

\vspace{0.1cm}
{\small \textcolor{errorstroke}{\textbf{Hardware constraint:}} AllReduce per layer requires 600--900 GB/s NVLink. \textbf{Intra-node only.}}

\end{frame}

\begin{frame}{The Jeff Dean Test}
\note{[2 min] A practical rule: if you attempt tensor parallelism across racks,
your training will stall. The communication volume per layer demands NVLink-class
bandwidth. For cross-rack scaling, use pipeline or data parallelism.}

\small
\begin{mlsyscard}{crimson}
\textbf{The Jeff Dean Test:} If you attempt \textbf{Tensor Parallelism} across server racks connected by standard Ethernet, your training will stall.
\end{mlsyscard}

\vspace{0.15cm}
\renewcommand{\arraystretch}{1.15}
{\footnotesize
\begin{tabular}{@{}llll@{}}
  \toprule
  \textbf{Strategy} & \textbf{Comm.\ per Step} & \textbf{BW Need} & \textbf{Placement} \\
  \midrule
  Tensor Par. & AllReduce per \emph{layer} & 600--900 GB/s & \textcolor{computestroke}{\textbf{Intra-node (NVLink)}} \\
  Pipeline Par. & P2P per \emph{stage boundary} & 25--50 GB/s & \textcolor{datastroke}{\textbf{Inter-node (InfiniBand)}} \\
  Data Par. & AllReduce per \emph{step} & 50--200 GB/s & \textcolor{routingstroke}{\textbf{Across replicas}} \\
  \bottomrule
\end{tabular}
}

\vspace{0.15cm}
\mlsysalert{Rule:}{Match communication frequency to interconnect bandwidth.}

\end{frame}

% =============================================================================
\section{Hybrid 3D}
% =============================================================================

\begin{frame}{Hybrid 3D Parallelism}
\note{[3 min] This is how frontier models are actually trained. Walk through the
example: GPT-4 scale on 8,192 H100s. TP=8 (within node), PP=16 (across node groups),
DP=64 (replicas). Total = 8,192 GPUs. The key is matching each parallelism dimension
to the right interconnect tier.}

% --- Full-width image ---
\centering
\safeimg[width=0.92\textwidth,height=4.8cm,keepaspectratio]{hybrid-3d.pdf}

\vspace{0.1cm}
{\small \textbf{Total} = $t \times p \times d = 8 \times 16 \times 64 = 8{,}192$ GPUs. Frequency: TP (kHz) $>$ PP (MHz) $>$ DP (Hz).}

\end{frame}

\begin{frame}{Parallelism Decision Tree}
\note{[2 min] Simplified but useful. First question: does the model fit in one GPU?
If yes, data parallelism. If no, does it fit in one node? If yes, TP+DP. If no,
full 3D parallelism. The decision is driven by constraints, not preference.}

\small
\begin{columns}[T]
  \begin{column}{0.55\textwidth}
    \textbf{Constraint Satisfaction:}

    \vspace{0.15cm}
    \begin{enumerate}\setlength\itemsep{2pt}
      \item {\footnotesize Model fits in 1 GPU?\\$\to$ \textcolor{computestroke}{\textbf{Data Parallelism}}}
      \item {\footnotesize Model fits in 1 node (8 GPUs)?\\$\to$ \textcolor{computestroke}{\textbf{TP}} within node + \textcolor{routingstroke}{\textbf{DP}} across}
      \item {\footnotesize Model exceeds 1 node?\\$\to$ \textcolor{datastroke}{\textbf{PP}} across nodes + \textcolor{computestroke}{\textbf{TP}} + \textcolor{routingstroke}{\textbf{DP}}}
      \item {\footnotesize ZeRO/FSDP can extend each regime before escalating}
    \end{enumerate}
  \end{column}
  \begin{column}{0.42\textwidth}
    \begin{mlsyscard}{crimson}
    \textbf{Production Reality}\\[0.1cm]
    {\footnotesize
    \textbf{Llama 3 (405B)}: TP=8, PP=16, DP=varies\\
    \textbf{PaLM (540B)}: 6,144 TPU v4\\
    \textbf{GPT-4}: estimated 25K+ A100s\\[0.1cm]
    MFU: 40--46\% (best reported)\\
    \textcolor{errorstroke}{$>$50\% of peak is never achieved.}
    }
    \end{mlsyscard}
  \end{column}
\end{columns}

\end{frame}

% --- ACTIVE LEARNING 3: Discussion ---
\begin{frame}{Discussion: Design a Training Config}
\note{[3 min] Turn-and-talk. Students design a 3D config for a 70B model
on 64 A100s. No single right answer. The point is reasoning about constraints.
Cold-call 2--3 pairs. Possible answers: TP=8, PP=2, DP=4 or TP=8, PP=4, DP=2.}

\centering
\vspace{0.8cm}
{\large\bfseries Turn and Talk \textcolor{midgray}{(90 seconds)}}

\vspace{0.5cm}
{\large You have \textbf{64 A100 GPUs} (80 GB each) across \textbf{8 nodes}.\\
You must train a \textbf{70B parameter} model.\\[0.3cm]
\alert{Design a 3D parallelism configuration: $(t, p, d)$ = ?}}

\vspace{0.4cm}
\begin{columns}[c]
  \begin{column}{0.25\textwidth}\centering\small Tensor (t)\\ {\scriptsize within node}\end{column}
  \begin{column}{0.25\textwidth}\centering\small Pipeline (p)\\ {\scriptsize across nodes}\end{column}
  \begin{column}{0.25\textwidth}\centering\small Data (d)\\ {\scriptsize replicas}\end{column}
  \begin{column}{0.15\textwidth}\centering\small Total\\ {\scriptsize = 64?}\end{column}
\end{columns}

\end{frame}

% =============================================================================
\section{RLHF at Scale}
% =============================================================================

\begin{frame}{Multi-Model Training: RLHF}
\note{[3 min] RLHF breaks the single-model assumption. PPO requires 4 models:
policy (training), reference (inference), reward (inference), value (training).
The aggregate memory dwarfs standard pre-training. DPO simplifies to 2 models,
halving GPU requirements.}

\small
\begin{columns}[T]
  \begin{column}{0.55\textwidth}
    \textbf{PPO requires 4 concurrent models:}
    \vspace{0.1cm}
    {\footnotesize
    \renewcommand{\arraystretch}{1.15}
    \begin{tabular}{@{}llr@{}}
      \toprule
      \textbf{Model} & \textbf{Mode} & \textbf{Memory} \\
      \midrule
      Policy (70B)   & Training  & 1,120 GB \\
      Reference (70B) & Inference & 140 GB \\
      Reward (13B)   & Inference & 26 GB \\
      Value (13B)    & Training  & 208 GB \\
      \midrule
      \textbf{Total} & & \textbf{1,494 GB} \\
      \bottomrule
    \end{tabular}
    }
  \end{column}
  \begin{column}{0.42\textwidth}
    \begin{mlsyscard}{datastroke}
    \textbf{DPO simplifies to 2 models}\\[0.1cm]
    {\footnotesize
    Policy: 1,120 GB (training)\\
    Reference: 140 GB (inference)\\
    \textbf{Total: 1,260 GB}\\[0.1cm]
    48\% less memory\\
    No generation phase\\
    2$\times$ higher MFU (40--55\%)
    }
    \end{mlsyscard}
  \end{column}
\end{columns}

\end{frame}

% =============================================================================
\section{Wrap-Up}
% =============================================================================

\begin{frame}{Fallacies}
\note{[2 min] Four common misconceptions. Each has quantitative evidence.}

\footnotesize
\textbf{Fallacy:} \textit{Linear speedup is achievable with sufficient engineering effort.}\\
A 1,000-GPU cluster will never train 1,000$\times$ faster. 100--200$\times$ is typical for communication-heavy workloads.

\vspace{0.08cm}
\textbf{Fallacy:} \textit{Data parallelism scales indefinitely by adding more GPUs.}\\
Beyond the critical batch size ($\sim$8K ResNet, $\sim$4M tokens LLMs), doubling GPUs doubles cost without convergence speedup.

\vspace{0.08cm}
\textbf{Fallacy:} \textit{FSDP and ZeRO always improve training efficiency.}\\
7B model on 80 GB GPUs: DDP = 145 samples/s vs.\ FSDP = 118 samples/s (19\% slower) when memory is not the bottleneck.

\vspace{0.08cm}
\textbf{Fallacy:} \textit{Parallelism overhead is constant regardless of model size.}\\
1B model: AllReduce = 33\% of step. 70B model: AllReduce = 7\%. Decisions invert at scale.

\end{frame}

\begin{frame}{Pitfalls}
\note{[2 min] Three operational pitfalls. If short: cover just the first two.}

\footnotesize
\textbf{Pitfall:} \textit{Hyperparameters tuned on small clusters transfer to large-scale training.}\\
Learning rate must scale with batch size: linear rule ($\eta \propto k$) with warmup, or $\sqrt{k}$ above critical batch size. A team training ResNet-50 at batch 32K with linear scaling got 74.8\% vs.\ 76.2\% at batch 256.

\vspace{0.12cm}
\textbf{Pitfall:} \textit{Choosing parallelism strategy based solely on memory constraints.}\\
For 175B on 64 A100s: pipeline parallelism with 8 stages achieves 23\% higher throughput than tensor parallelism degree-8, despite similar memory footprints. Profile communication patterns, not just memory.

\vspace{0.12cm}
\textbf{Pitfall:} \textit{Gradient accumulation is free.}\\
Accumulated FP16 gradients risk overflow when summing hundreds of tensors. A team using 16-step accumulation in FP16 experienced loss spikes at step 1,200; 4-step accumulation resolved the instability.

\end{frame}

% --- RETRIEVAL PRACTICE ---

% --- MUDDIEST POINT ---
\begin{frame}{Muddiest Point}
\note{[2 min] Quick anonymous poll. Students write on a slip of paper or submit
digitally. Collect and scan for patterns. Address the top 2--3 confusions in the
next lecture's opening. This closes the feedback loop.}

\centering
\vspace{1.0cm}
{\Large\bfseries What was the \alert{muddiest point} today?}

\vspace{0.8cm}
{\normalsize Write down the concept you found \textbf{most confusing}.}

\vspace{0.5cm}
{\small\textcolor{midgray}{Anonymous. One sentence. Submit before you leave.}}

\end{frame}

\begin{frame}{What Were the Key Ideas?}
\note{[2 min] Retrieval practice. Students write 90 seconds, no notes.
Walk around the room.}

\centering
\vspace{1.5cm}
{\Large\bfseries Close your notes.}

\vspace{0.8cm}
{\large Write down the \textbf{4 most important concepts} from today.}

\vspace{0.8cm}
{\normalsize\textcolor{midgray}{90 seconds --- no peeking.}}

\end{frame}

\begin{frame}{Key Takeaways}
\note{[2 min] Reveal. Walk through each bullet. Emphasize quantitative anchors:
355 years single-GPU, 18\% bubble tax, 50\% MFU ceiling, Nx ZeRO savings.}

\scriptsize
\begin{itemize}\setlength\itemsep{0pt}
  \item \textbf{Communication-Computation Ratio}: Determines whether your cluster is a supercomputer or idling heaters.
  \item \textbf{Data parallelism}: Default. Scales linearly until the \textbf{Critical Batch Size} (8K--4M tokens).
  \item \textbf{ZeRO/FSDP}: Shard optimizer (4$\times$), gradients (8$\times$), parameters ($N\times$). Cost: 10--25\% comm overhead.
  \item \textbf{Tensor parallelism}: Splits matrix ops within layers. Requires 600--900 GB/s NVLink. Intra-node only.
  \item \textbf{Pipeline parallelism}: Splits layers across stages. Low BW need. Bubble: $(P\!-\!1)/(P\!-\!1\!+\!M)$.
  \item \textbf{Hybrid 3D}: TP (node) + PP (cross-node) + DP (replicas). Total = $t \times p \times d$.
  \item \textbf{Linear Scaling Rule}: LR $\propto k$ when batch $\propto k$, with warmup. Beyond $B^*$, use $\sqrt{k}$.
\end{itemize}

\end{frame}

\begin{frame}{References}
\note{[1 min] Point students to canonical papers.}

\small
\mlsysref{Goyal+17}{Goyal et al. ``Accurate, Large Minibatch SGD.'' arXiv:1706.02677, 2017.}
\mlsysref{Rajbhandari+20}{Rajbhandari et al. ``ZeRO: Memory Optimization Towards Training Trillion Parameter Models.'' SC 2020.}
\mlsysref{Shoeybi+20}{Shoeybi et al. ``Megatron-LM: Training Multi-Billion Parameter Language Models.'' arXiv:1909.08053, 2020.}
\mlsysref{Narayanan+21}{Narayanan et al. ``Efficient Large-Scale Language Model Training on GPU Clusters.'' SC 2021.}
\mlsysref{Zhao+23}{Zhao et al. ``PyTorch FSDP: Experiences on Scaling Fully Sharded Data Parallel.'' VLDB 2023.}
\mlsysref{Llama3+24}{Meta AI. ``The Llama 3 Herd of Models.'' arXiv:2407.21783, 2024.}

\end{frame}

\begin{frame}{Next Lecture: Collective Communication}
\note{[1 min] Forward hook. We defined the logical traffic patterns. Next chapter
opens the black box of AllReduce, AllGather, ReduceScatter --- the actual algorithms
that move bytes through rings, trees, and rails.}

\footnotesize
\begin{columns}[c]
  \begin{column}{0.30\textwidth}
    \centering
    \begin{mlsyscard}{computestroke}
    {\large\bfseries Ring AllReduce}\\[0.1cm]
    {\scriptsize Bandwidth-optimal\\$2(N\!-\!1)/N$ utilization}
    \end{mlsyscard}
  \end{column}
  \begin{column}{0.30\textwidth}
    \centering
    \begin{mlsyscard}{datastroke}
    {\large\bfseries Tree AllReduce}\\[0.1cm]
    {\scriptsize Latency-optimal\\$O(\log N)$ steps}
    \end{mlsyscard}
  \end{column}
  \begin{column}{0.30\textwidth}
    \centering
    \begin{mlsyscard}{routingstroke}
    {\large\bfseries Hierarchical}\\[0.1cm]
    {\scriptsize NVLink intra-node\\InfiniBand inter-node}
    \end{mlsyscard}
  \end{column}
\end{columns}

\vspace{0.2cm}
\centering
{\normalsize From the \emph{logic} of what to send to the \emph{mechanics} of how to route it.}\\[0.1cm]
{\small\textbf{The algorithms that make AllReduce, AllGather, and ReduceScatter efficient at scale.}}

\end{frame}



\appendix

\begin{frame}{Backup: ZeRO Memory Calculation Reference}
\note{Backup: Extended ZeRO memory calculation.}

\footnotesize
\textbf{Memory per GPU with ZeRO stages (N GPUs, M parameters):}

\vspace{0.15cm}
\renewcommand{\arraystretch}{1.2}
\begin{tabular}{@{}lll@{}}
  \toprule
  \textbf{Stage} & \textbf{Memory/GPU} & \textbf{7B on 64 GPUs} \\
  \midrule
  DDP (baseline) & $2M + 2M + 12M = 16M$ bytes & 112 GB \\
  ZeRO-1 (opt)   & $2M + 2M + 12M/N$ & 29.3 GB \\
  ZeRO-2 (+grad) & $2M + (2M + 12M)/N$ & 15.5 GB \\
  ZeRO-3 (+param)& $(2M + 2M + 12M)/N$ & 1.75 GB \\
  \bottomrule
\end{tabular}

\vspace{0.1cm}
{\footnotesize ZeRO-3 trades 10--25\% communication overhead for $N\times$ memory reduction.}

\end{frame}

\begin{frame}{Backup: Pipeline Schedule Comparison}
\note{Backup: GPipe vs 1F1B comparison.}

\footnotesize
\renewcommand{\arraystretch}{1.15}
\begin{tabular}{@{}lllr@{}}
  \toprule
  \textbf{Schedule} & \textbf{Bubble} & \textbf{Peak Memory} & \textbf{Best For} \\
  \midrule
  GPipe    & $(P-1)/(P-1+M)$ & $O(M \times P)$ & Simplicity \\
  1F1B     & $(P-1)/(P-1+M)$ & $O(P)$          & Memory-constrained \\
  Interleaved & Reduced         & $O(P)$          & Many stages \\
  \bottomrule
\end{tabular}

\vspace{0.15cm}
{\footnotesize 1F1B interleaves forward and backward passes to limit peak memory to $O(P)$ instead of $O(M \times P)$.}

\end{frame}


\end{document}
